%%%%%%%%%%%%%%%%%%%%%%%%%%%%%%%%%%%%%%%%%%%%%%%%%%%%%%%%%%%%%%%%%%%%%%%%%%%
% Función: capítulo de ejemplo de conclusiones de la tesis por compendio.
%
% Uso: copia de restauración; normalmente no debes modificarla.
% Tras restaurarla, modifica la copia activa para adaptarla a tu trabajo.
%%%%%%%%%%%%%%%%%%%%%%%%%%%%%%%%%%%%%%%%%%%%%%%%%%%%%%%%%%%%%%%%%%%%%%%%%%%

\chapter{Conclusiones y líneas futuras}
\label{cha:compendium-conclusions}

Como ejemplo de conexión entre conclusiones y alcance, compara los objetivos ligados a la aplicación en evaluación funcional~\cite{nieva2024proposal}, el contexto del sistema descrito para EYEFUL~\cite{nieva2023eyeful} y el escenario de sensado diferente de la localización de barcos~\cite{ramirez2026jstars}. Utiliza la comparación para preguntarte qué afirmaciones están sustentadas por tus propias evidencias y cuáles exigirían más validación. No presentes resultados de esas publicaciones como resultados de esta tesis ficticia.

\section{Conclusiones principales}

Resume las conclusiones sustentadas por el conjunto de la tesis, relaciónalas con los objetivos y discute las principales limitaciones. Cierra el resumen extendido identificando las líneas de investigación que se derivan de los resultados.

\section{Limitaciones}

Discute las hipótesis y restricciones prácticas que limitan la generalización de los resultados y deben considerarse al interpretar las conclusiones.

\section{Líneas futuras de investigación}
\subsection{Extensiones inmediatas}

Identifica las mejoras y experimentos que se derivan directamente del trabajo presentado en la tesis.

\subsection{Oportunidades de investigación a largo plazo}

Plantea cuestiones más amplias abiertas por el conjunto de las contribuciones y publicaciones.

%%% Local Variables:
%%% TeX-master: "../../book"
%%% End:
